% Capítulo sin numeración
\chapter*{Resumen}

El sistema operativo se encarga de gestionar los recursos de hardware de un dispositivo electrónico y de brindar servicios a los programas de aplicación. La gestión de los procesos e hilos y los recursos de un sistema operativo son fundamentales para su correcto funcionamiento y desempeño. En este contexto, la planificación es una tarea crítica para asegurar el uso eficiente de los recursos del sistema.

Este trabajo parte de un proyecto integrador previo desarrollado por dos alumnos de la FCEFyN, que realizaron un planificador de corto plazo modelado mediante Redes de Petri. Las Redes de Petri son una herramienta matemática formal de modelado que permite representar y analizar sistemas de eventos discretos, como es el caso de la planificación de procesos. En este contexto, permiten describir con claridad los estados y eventos del planificador, y ofrecen una base adecuada para incorporar nuevas reglas de asignación de recursos.

En el marco de esta investigación, se adaptó el proyecto integrador previo desde FreeBSD 11 hasta FreeBSD 13.1. Esta actualización permitió trabajar sobre una base de código más reciente que la del desarrollo original y dejar el planificador en condiciones de ser extendido sobre esa versión.

Luego de la actualización, se incorporaron dos módulos al planificador. El primero, denominado en este informe módulo de inhibición de asignación de trabajo a procesadores, permite impedir manualmente que determinados núcleos reciban nuevos hilos en sus colas de ejecución, sin apagar físicamente el procesador ni intervenir sobre la administración energética del hardware. El segundo, denominado módulo de monopolización de núcleos, permite reservar un núcleo para la ejecución de un hilo específico. Ambas funcionalidades se presentan como mecanismos de control del planificador y como base para futuros trabajos que incorporen políticas automáticas, mediciones de rendimiento o estrategias de eficiencia energética.
